% =====================================================
% 题型：三角函数图象与性质综合解答题 (q_trig_graph_properties.tex)
% =====================================================
% =====================================================
% 难度：★★ (中档)
% =====================================================
\newcommand{\qTrigGraphPropertiesStem}{已知函数 \(f(x) = A\sin(\omega x + \varphi)\)（其中 \(x \in \mathbb{R}\)，\(A > 0\)，\(\omega > 0\)，\(|\varphi| < \dfrac{\pi}{2}\)）的部分图象如图所示。\\
（1）求 \(f(x)\) 的解析式；\\
（2）求 \(f(x)\) 的单调递增区间及 \(f(x)\) 在区间 \(\left[0, \dfrac{\pi}{2}\right]\) 上的最大值与最小值。}

\newcommand{\qTrigGraphPropertiesAnswer}{（1）\(f(x) = 2\sin\left(2x - \dfrac{\pi}{3}\right)\)；（2）递增区间为 \(\left[-\dfrac{\pi}{12} + k\pi, \dfrac{5\pi}{12} + k\pi\right]\) (\(k \in \mathbb{Z}\))，在 \([0, \pi/2]\) 上最大值为 \(2\)，最小值为 \(-\sqrt{3}\)}

\newcommand{\qTrigGraphPropertiesFigure}[1][1.0]{%
  \begin{tikzpicture}[scale=#1, >=stealth, x=1.5cm, y=0.8cm]
    % Draw axes
    \draw[->, thick] (-0.8, 0) -- (3.5, 0) node[right] {\small \(x\)};
    \draw[->, thick] (0, -2.5) -- (0, 2.5) node[above] {\small \(y\)};
    \node[below left] at (0,0) {\small \(O\)};
    
    % Plot function y = 2*sin(2*x - pi/3)
    % Domain: -0.2 to 3.3
    \draw[very thick, domain=-0.2:3.3, samples=100] plot (\x, {2*sin(deg(2*\x - 3.14159/3))});
    
    % Dotted lines for amplitude
    \draw[dashed, gray] (-0.8, 2) -- (3.3, 2);
    \draw[dashed, gray] (-0.8, -2) -- (3.3, -2);
    \node[left] at (0, 2) {\small \(2\)};
    \node[left] at (0, -2) {\small \(-2\)};
    
    % Key points
    \coordinate (Z1) at (0.5236, 0);   % pi/6
    \coordinate (Z2) at (2.0944, 0);   % 2pi/3
    
    \fill (Z1) circle (1.5pt) node[below left, yshift=2pt] {\small \(\dfrac{\pi}{6}\)};
    \fill (Z2) circle (1.5pt) node[below right, yshift=2pt] {\small \(\dfrac{2\pi}{3}\)};
  \end{tikzpicture}%
}

\newcommand{\qTrigGraphPropertiesSolution}{%
  【标准解答】
  
  （1）由函数图象可知，函数的最大值为 \(2\)，且 \(A > 0\)，
  所以：
  \[
    A = 2.
  \]
  图象与 \(x\) 轴在正半轴的相邻两个零点分别为 \(x_1 = \dfrac{\pi}{6}\) 和 \(x_2 = \dfrac{2\pi}{3}\)。
  两相邻零点之间的距离为半个周期，即：
  \[
    \frac{T}{2} = x_2 - x_1 = \frac{2\pi}{3} - \frac{\pi}{6} = \frac{\pi}{2} \implies T = \pi.
  \]
  因为 \(T = \dfrac{2\pi}{\omega}\)，且 \(\omega > 0\)，所以：
  \[
    \frac{2\pi}{\omega} = \pi \implies \omega = 2.
  \]
  于是，函数的解析式可表述为：
  \[
    f(x) = 2\sin(2x + \varphi).
  \]
  根据图象，当 \(x = \dfrac{\pi}{6}\) 时，函数图象处于上升沿穿过零点（即正弦函数的最基本零点 \(0\)），
  因此满足：
  \[
    2 \cdot \frac{\pi}{6} + \varphi = 0 \implies \frac{\pi}{3} + \varphi = 0 \implies \varphi = -\frac{\pi}{3}.
  \]
  因为 \(|\varphi| < \dfrac{\pi}{2}\)，且 \(-\dfrac{\pi}{3}\) 在该范围内，符合题意。
  所以，所求函数的解析式为：
  \[
    f(x) = 2\sin\left(2x - \dfrac{\pi}{3}\right).
  \]
  
  \vspace{0.5em}
  
  （2）由正弦函数的单调性，令：
  \[
    -\frac{\pi}{2} + 2k\pi \le 2x - \dfrac{\pi}{3} \le \frac{\pi}{2} + 2k\pi, \quad k \in \mathbb{Z}.
  \]
  各部分同加 \(\dfrac{\pi}{3}\) 得：
  \[
    -\frac{\pi}{6} + 2k\pi \le 2x \le \frac{5\pi}{6} + 2k\pi, \quad k \in \mathbb{Z}.
  \]
  同除以 \(2\) 得：
  \[
    -\frac{\pi}{12} + k\pi \le x \le \frac{5\pi}{12} + k\pi, \quad k \in \mathbb{Z}.
  \]
  所以，\(f(x)\) 的单调递增区间为：
  \[
    \left[-\frac{\pi}{12} + k\pi, \frac{5\pi}{12} + k\pi\right] \quad (k \in \mathbb{Z}).
  \]
  
  \vspace{0.5em}
  
  当 \(x \in \left[0, \dfrac{\pi}{2}\right]\) 时，有：
  \[
    2x \in [0, \pi] \implies 2x - \frac{\pi}{3} \in \left[-\frac{\pi}{3}, \frac{2\pi}{3}\right].
  \]
  设 \(t = 2x - \dfrac{\pi}{3}\)，则 \(t \in \left[-\dfrac{\pi}{3}, \dfrac{2\pi}{3}\right]\)。
  
  在正弦函数 \(y = \sin t\) 在此区间上的表现中：
  \begin{itemize}
    \item 当 \(t = \dfrac{\pi}{2}\)，即 \(2x - \dfrac{\pi}{3} = \dfrac{\pi}{2} \implies x = \dfrac{5\pi}{12} \in \left[0, \dfrac{\pi}{2}\right]\) 时，\(\sin t\) 取得最大值 \(1\)，
    此时 \(f(x)\) 取得最大值：
    \[
      f_{\max} = 2 \cdot 1 = 2.
    \]
    \item 当 \(t = -\dfrac{\pi}{3}\)，即 \(x = 0\) 时，\(\sin t\) 取得最小值 \(\sin\left(-\dfrac{\pi}{3}\right) = -\dfrac{\sqrt{3}}{2}\)，
    此时 \(f(x)\) 取得最小值：
    \[
      f_{\min} = 2 \cdot \left(-\frac{\sqrt{3}}{2}\right) = -\sqrt{3}.
    \]
  \end{itemize}
  所以，\(f(x)\) 在 \(\left[0, \dfrac{\pi}{2}\right]\) 上的最大值为 \(2\)，最小值为 \(-\sqrt{3}\)。
}

% =====================================================
% 别名兼容定义
% =====================================================
\newcommand{\qTJWTrigGraphPropertiesStem}{\qTrigGraphPropertiesStem}
\newcommand{\qTJWTrigGraphPropertiesAnswer}{\qTrigGraphPropertiesAnswer}
\newcommand{\qTJWTrigGraphPropertiesFigure}{\qTrigGraphPropertiesFigure}
\newcommand{\qTJWTrigGraphPropertiesSolution}{\qTrigGraphPropertiesSolution}
